\section{Direct-sum separation and prepared-family bookkeeping}
\label{sec:app_bnt_separation_supplier}

We now prove the supporting separation and regrouping results used by the
source-facing block-injectivity and supplier theorems in
Chapter~\ref{ch:bnt}.  The direct-sum argument follows the ground-space
uniqueness mechanism of \cite{PerezGarcia2007Matrix}: disjoint local
MPS spaces give pair-trace selectors, products of those selectors isolate one
BNT block, and an injective prefix supplies an arbitrary matrix on that block.
The remaining results identify dimensions while phase-equivalent prepared
blocks are grouped into sectors and normalize the resulting finite weight
family.

\begin{lemma}[Three-block separation of distinct BNT representatives]
    \label{lem:bnt_direct_sum_injective_prefix_word_span}
    \lean{MPSTensor.groundSpace_inf_eq_bot_of_blocksNotGaugePhaseEquiv_same_dim_of_dim_ge_c1}
    \lean{MPSTensor.pairTraceSeparatingAt_threeBlock_of_blocksNotGaugePhaseEquiv_same_dim_of_dim_ge_c1}
    \lean{MPSTensor.forall_pairTraceSeparatingAt_threeBlock_of_blocksNotGaugePhaseEquiv_c1}
    \leanok
    \uses{def:bnt_cpsv, def:blocks_not_gauge_phase_equiv,
        def:left_canonical_tensor, def:mpv_overlap,
        thm:not_gauge_phase_eventually_not_proportional}
    Let $A^1,\ldots,A^r$ be irreducible left-canonical representatives with
    normalized self-overlaps, and suppose that distinct same-dimensional
    representatives are not gauge-phase equivalent. Fix $L_0>0$ such that
    every block is injective at lengths $L_0$, $L_0+1$, and $3(L_0+1)$.
    Then each ordered pair $k\ne j$ is trace-separated at the common length
    $S=3(L_0+1)$: the map
    \begin{align}
        (c_\omega)_{|\omega|=S}
         & \longmapsto
        \left(\sum_{|\omega|=S}c_\omega A^k_\omega,
        \sum_{|\omega|=S}c_\omega A^j_\omega\right)
        \notag
    \end{align}
    is onto $\MN{D_k}\oplus\MN{D_j}$. In particular, one may prescribe an
    arbitrary matrix on either block while making the same word combination
    vanish on the other block.
\end{lemma}

\begin{proof}
    \leanok
    \uses{thm:not_gauge_phase_eventually_not_proportional}
    Put $L=L_0+1$ and $S=3L$. For a block $A$, write
    \begin{align}
        \Gamma_n^A(X)
         & =\sum_{|w|=n}\tr(XA_w)\ket{w},
         &
        \mathcal G_n(A)
         & =\operatorname{im}\Gamma_n^A.
        \notag
    \end{align}
    Block injectivity at $L$ and $S$ makes both boundary maps injective; in
    particular,
    \begin{align}
        \dim\mathcal G_L(A^q) & =D_q^2
        \qquad (1\leq q\leq r).
        \label{eq:app_bnt_local_space_dimension}
    \end{align}

    Fix $k\ne j$ and first order the pair so that $D_j\leq D_k$. Suppose that
    a nonzero vector belongs to
    $\mathcal G_S(A^k)\cap\mathcal G_S(A^j)$. There are boundary matrices
    $X$ and $Y$ such that
    \begin{align}
        0\ne\psi
         & =\Gamma_S^{A^k}(X)=\Gamma_S^{A^j}(Y).
        \label{eq:app_bnt_common_three_block_vector}
    \end{align}
    Injectivity at $S$ gives $X\ne0$ and $Y\ne0$. Split every length-$S$ word
    as $uvw$ with $|u|=|v|=|w|=L$. Equality
    \ref{eq:app_bnt_common_three_block_vector} gives
    \begin{align}
        \tr(XA^k_uA^k_vA^k_w)
         & =\tr(YA^j_uA^j_vA^j_w).
        \label{eq:app_bnt_three_block_trace_relation}
    \end{align}
    Because the length-$L$ words of $A^k$ span $\MN{D_k}$ and $X\ne0$, the
    two-sided ideal generated by $X$ is all of $\MN{D_k}$:
    \begin{align}
        \spn\{A^k_uXA^k_v:|u|=|v|=L\}
         & =\MN{D_k}.
        \label{eq:app_bnt_two_sided_span}
    \end{align}
    Indeed, expand arbitrary left and right matrix units in length-$L$ words;
    a nonzero entry of $X$ then produces every matrix unit. Thus, for each
    $Z\in\MN{D_k}$, choose coefficients $a_{uv}$ with
    $Z=\sum_{u,v}a_{uv}A^k_uXA^k_v$ and set
    $R(Z)=\sum_{u,v}a_{uv}A^j_uYA^j_v$. Cyclicity of the trace and
    \ref{eq:app_bnt_three_block_trace_relation}, applied to the word $vwu$,
    yield
    \begin{align}
        \tr(ZA^k_w)
         & =\sum_{u,v}a_{uv}\tr(XA^k_vA^k_wA^k_u)
        \notag                                    \\
         & =\sum_{u,v}a_{uv}\tr(YA^j_vA^j_wA^j_u)
        =\tr(R(Z)A^j_w).
        \notag
    \end{align}
    Hence $\Gamma_L^{A^k}(Z)=\Gamma_L^{A^j}(R(Z))$ for every $Z$, and therefore
    \begin{align}
        \mathcal G_L(A^k) & \subseteq\mathcal G_L(A^j).
        \label{eq:app_bnt_local_space_inclusion}
    \end{align}
    Equations~\ref{eq:app_bnt_local_space_dimension} and
    \ref{eq:app_bnt_local_space_inclusion} give
    \begin{align}
        D_k^2
         & =\dim\mathcal G_L(A^k)
        \leq\dim\mathcal G_L(A^j)=D_j^2.
        \notag
    \end{align}
    Together with $D_j\leq D_k$, this forces $D_j=D_k$ and equality in
    \ref{eq:app_bnt_local_space_inclusion}. This proves the unequal-dimension
    branch immediately: if $D_j<D_k$, a nonzero intersection is impossible.
    Reversing the initial ordering treats $D_k<D_j$ in the same way.

    It remains to exclude the intersection when $D_j=D_k$. Gauge-phase
    separation and Theorem~\ref{thm:not_gauge_phase_eventually_not_proportional}
    give, beyond any prescribed threshold, a chain length $N>L$ for which the
    two periodic MPVs are not proportional. On the other hand, equality
    $\mathcal G_L(A^k)=\mathcal G_L(A^j)$ gives identical translated
    length-$L$ constraints on the $N$-site ring. The parent-Hamiltonian
    uniqueness mechanism identifies the corresponding chain ground spaces
    with the one-dimensional MPV spans:
    \begin{align}
        \mathcal G_{N,L}(A^k)
         & =\C\,V^{(N)}(A^k),
         &
        \mathcal G_{N,L}(A^j)
         & =\C\,V^{(N)}(A^j).
        \notag
    \end{align}
    Equality of the local spaces therefore forces these two MPV lines to be
    equal, contradicting the choice of $N$. Consequently, for every distinct
    pair, including the unequal-dimensional pairs,
    \begin{align}
        \mathcal G_S(A^k)\cap\mathcal G_S(A^j) & =\{0\}.
        \label{eq:app_bnt_three_block_intersection_zero}
    \end{align}

    Finally, let $\Phi_{kj}$ be the displayed length-$S$ word-evaluation map
    from coefficient families to $\MN{D_k}\oplus\MN{D_j}$. Pair the target by
    \begin{align}
        \langle (X,Y),(M,N)\rangle
         & =\tr(XM)+\tr(YN).
        \notag
    \end{align}
    A pair $(X,Y)$ annihilates $\operatorname{im}(\Phi_{kj})$ precisely when
    \begin{align}
        \tr(XA^k_w)+\tr(YA^j_w) & =0
        \qquad (|w|=S),
        \notag
    \end{align}
    or equivalently
    $\Gamma_S^{A^k}(X)=\Gamma_S^{A^j}(-Y)$. Via the injective boundary maps,
    this is the canonical identification
    \begin{align}
        \operatorname{im}(\Phi_{kj})^\perp
         & =\mathcal G_S(A^k)\cap\mathcal G_S(A^j).
        \label{eq:app_bnt_pair_map_annihilator}
    \end{align}
    The right-hand side vanishes by
    \ref{eq:app_bnt_three_block_intersection_zero}. The trace pairing is
    nondegenerate and all spaces are finite-dimensional, so
    $\operatorname{im}(\Phi_{kj})^\perp=\{0\}$ implies that $\Phi_{kj}$ is
    surjective. Surjectivity first gives an arbitrary value on one block and
    zero on an anchor block; the selector construction then chooses
    $(\Id,0)$ for the remaining pairs, and products of these homogeneous
    selectors give the simultaneous word-tuple span in
    Lemma~\ref{lem:bnt_direct_sum_block_separation_word_span}.
\end{proof}

\begin{lemma}[Common injective blocking for primitive irreducible families]%
    \label{thm:common_injective_blocking_tp_primitive_irreducible_family}
    \lean{MPSTensor.exists_common_injective_blocking_of_tp_primitive_irr_family}
    \leanok
    \uses{thm:tp_primitive_irred_block_injective,
        thm:tp_primitive_irreducible_extra_blocking}
    Let $(B_k)_{k=1}^r$ be a finite family of positive-bond-dimension MPS tensors,
    each trace-preserving, primitive, and tensor-irreducible. Then there is a
    single $L\ge 1$ such that each $B_k^{[L]}$ is one-site injective, and this
    blocking preserves trace preservation, primitivity, and irreducibility for
    every block.
\end{lemma}

\begin{proof}\leanok
    For each block,
    Theorem~\ref{thm:tp_primitive_irred_block_injective} gives a positive length
    at which it becomes one-site injective; the product of these lengths is a
    common length. Injectivity persists under positive further blocking, and
    Theorem~\ref{thm:tp_primitive_irreducible_extra_blocking} preserves the other
    hypotheses.
\end{proof}

\begin{lemma}[Total dimension after grouping by BNT representatives]%
    \label{lem:collapsed_bnt_sector_decomp_total_dim}
    \lean{MPSTensor.collapsedBntSectorDecomp_totalDim_eq_sum_dim}
    \leanok
    \uses{def:mpv_phase_class_data}
    Suppose CF summands are grouped by BNT representative, every original copy
    weight is nonzero, and every summand assigned to a representative has that
    representative's bond dimension. Then the total bond dimension of the
    resulting BNT sector decomposition is the sum of the original block
    dimensions.
\end{lemma}

\begin{proof}\leanok
    The collapsed decomposition counts one representative dimension for each
    assigned copy. Let $e(j,q)$ denote the original CF summand assigned to the
    $q$th copy of representative $P_j$. Since $(j,q)\mapsto e(j,q)$ reindexes
    the original summands and each assigned summand has the representative's bond
    dimension,
    \begin{align}
        \sum_{j=1}^{g}\sum_{q=1}^{r_j}\dim(P_j)
         & =
        \sum_{j=1}^{g}\sum_{q=1}^{r_j}D_{e(j,q)}
        =
        \sum_{k=1}^{r}D_k.
        \notag
    \end{align}
    This is the equality between the collapsed total dimension and the original
    direct-sum dimension.
\end{proof}

\begin{lemma}[Prepared phase-equivalent NTs have the same bond dimension]%
    \label{lem:prepared_phase_equiv_dim_eq}
    \lean{MPSTensor.dim_eq_of_MPVBlockPhaseEquiv_of_tp_primitive_irr}
    \leanok
    \uses{def:mpv_block_phase_equiv, thm:overlap_decay_rect_irr_tp}
    Let $X$ and $Y$ be trace-preserving, primitive, irreducible NTs of positive
    bond dimension. If their MPV families differ by a nonzero scalar power, then
    $X$ and $Y$ have the same bond dimension.
\end{lemma}

\begin{proof}\leanok
    \uses{thm:rect_overlap_norm_one_dim_eq}
    Choose $\zeta\neq0$ with
    $\ket{V^{(N)}(Y)}=\zeta^N\ket{V^{(N)}(X)}$ for every positive length~$N$.
    The normalized self-overlaps give $|O_{XX}(N)|\to 1$, and the phase relation
    gives
    \begin{align}
        O_{YX}(N)   & =\zeta^N O_{XX}(N),
        \notag                                \\
        |O_{YX}(N)| & =|\zeta|^N |O_{XX}(N)|.
        \notag
    \end{align}
    Since the same normalization forces $|\zeta|=1$, this gives
    $|O_{YX}(N)|\to 1$.
    If the two bond dimensions differed, the rectangular overlap-decay theorem
    would force $O_{YX}(N)\to 0$, contradicting the preceding limit.
\end{proof}

\begin{lemma}[Prepared BNT representative classes preserve bond dimension]%
    \label{lem:prepared_phase_class_dim_eq}
    \lean{MPSTensor.mpvPhaseClassData_dim_eq_of_tp_primitive_irr}
    \leanok
    \uses{def:mpv_phase_class_data, lem:prepared_phase_equiv_dim_eq}
    In a prepared family of positive-bond-dimension trace-preserving primitive
    irreducible NTs, every CF summand assigned to a BNT representative has that
    representative's bond dimension.
\end{lemma}

\begin{proof}\leanok
    Apply Lemma~\ref{lem:prepared_phase_equiv_dim_eq} to each representative and
    each assigned summand.
\end{proof}

\begin{lemma}[Prepared collapsed BNT total dimension]%
    \label{lem:prepared_collapsed_bnt_sector_decomp_total_dim}
    \lean{MPSTensor.collapsedBntSectorDecomp_totalDim_eq_sum_dim_of_tp_primitive_irr}
    \leanok
    % Source: canonical-form eq:II_CF1 and prop:char-BNT (Cirac2016MPDO_arXiv),
    % with the equalMPS dimension conclusion (Cirac2016MPDO_arXiv).
    For prepared trace-preserving primitive irreducible NTs of positive bond
    dimension with all copy weights nonzero, grouping CF summands by BNT
    representative preserves the total bond dimension: the collapsed sector
    decomposition has total bond dimension equal to the sum of the original block
    dimensions.
\end{lemma}

\begin{proof}\leanok\uses{lem:collapsed_bnt_sector_decomp_total_dim,
        lem:prepared_phase_class_dim_eq}
    Every assigned summand shares the bond dimension of its representative, by
    the equal-dimension fact above. The displayed regrouping identity in
    Lemma~\ref{lem:collapsed_bnt_sector_decomp_total_dim} then identifies the
    collapsed sum over pairs $(j,q)$ with $\sum_{k=1}^{r}D_k$, the original
    direct-sum dimension.
\end{proof}

\begin{theorem}[Conditional BNT form after blocking for an arbitrary tensor]%
    \label{thm:paperbnt_supplier_after_blocking}
    \lean{MPSTensor.exists_isBNTCanonicalForm_afterBlocking_pos}
    \leanok
    \uses{thm:prepared_bnt_blocks_after_blocking,
        thm:paperbnt_supplier_prepared,
        def:same_mpv2_pos}
    Let $A$ be an MPS tensor. Then there exists a blocking length $p \ge 1$,
    nonzero weights $\mu_k$, and blocks $B_k$ over the blocked alphabet, such
    that:
    \begin{enumerate}
        \item every block has positive bond dimension;
        \item every block is TP-normalized, primitive, irreducible, and
              injective;
        \item $A^{[p]}$ and $\bigoplus_k \mu_k B_k$ generate the same MPV family
              at every positive length.
    \end{enumerate}
    Moreover, under $|\mu_k| \le 1$ for every $k$ and $|\mu_k|=1$ for at least
    one $k$, there exists a sector decomposition $P$ over the blocked alphabet in
    the BNT canonical form of Definition~\ref{def:sector_bnt_cf} such that
    $A^{[p]}$ and $P$ generate the same MPV family at every positive length.
\end{theorem}

\begin{proof}\leanok
    The after-blocking reduction supplies the TP-normalized primitive irreducible
    injective blocks and the positive-length MPV identity. The normalized-weight
    hypotheses are the choice that
    \cite[Section~2.3]{Cirac2016MPDO_arXiv} makes;
    Theorem~\ref{thm:paperbnt_supplier_prepared} then produces the sector
    decomposition, following the BNT selection of
    \cite[Appendix~A]{Cirac2016MPDO_arXiv}. The remaining rescaling step is
    recorded in Remark~\ref{rem:arbitrary_input_reduction_gap}.
\end{proof}

\begin{theorem}[Projective weight normalization]%
    \label{thm:paperbnt_weight_normalization}
    \lean{MPSTensor.exists_weight_normalization}
    \leanok
    For nonzero weights $\mu_1,\ldots,\mu_r$ with $r\geq1$ there is a positive
    real $m$ such that every normalized weight satisfies
    \begin{align}
        \Bigl|\frac{\mu_k}{m}\Bigr| & \leq1,
        \notag                               \\
        \frac{\mu_k}{m}             & \neq0,
        \notag
    \end{align}
    and at least one normalized weight has unit modulus,
    $|\mu_{k_0}/m|=1$.
    This is the normalization that
    \cite[line~246]{Cirac2016MPDO_arXiv} chooses for the canonical-form
    weights.
\end{theorem}

\begin{proof}\leanok
    Take $m=\max_k|\mu_k|$; the maximum over a nonempty finite
    family is attained, positive because every weight is nonzero, an upper
    bound for every modulus, and equal to the modulus of the attaining
    weight.
\end{proof}

\begin{lemma}[Common scalar rescaling of block weights]%
    \label{lem:mpv_to_tensor_from_blocks_weight_mul_left}
    \lean{MPSTensor.mpv_toTensorFromBlocks_weight_mul_left}
    \leanok
    Multiplying every block weight in a block-diagonal MPS tensor by a scalar
    $c$ multiplies every length-$N$ MPV coefficient by $c^N$.
\end{lemma}

\begin{proof}\leanok
    \uses{thm:mpv_decomposition}
    Expand the block-diagonal MPV coefficient as a finite sum over blocks and
    factor $c^N$ out of the sum.
\end{proof}
